\section{Conclusion}\label{sec:conclusion}

If local summaries preserve the evidence a task oracle needs, and local merges preserve the cross-span interactions, then preference optimization can be carried out on compression trees rather than full documents. The theorem ladder makes this precise: the structural bundle gives root preservation, the optimization bundle gives objective equivalence for DPO/GRPO, and the certificate bundle turns sampled audits into a finite-sample bound on preference drift.

The Markov changepoint experiments show the mechanism working in practice: decomposing documents into trees of subproblems matches or improves on a flat baseline, the supervision itself can be distributed across the tree, and both results hold across two DGPs and the full range of supervision budgets. The same structural mechanism---preserving cross-span interactions through learned merges---applies to semantic tasks such as manifesto RILE scoring, where the sufficient statistic is a high-dimensional representation and the merge rule is a neural network. The Semantic Forests companion extends these ideas to orchestration and deployment at system scale.
